\begin{definition}[Wild quadratic coefficient table]\label{mod:cases-wildquadratic-coefficienttable}
Put \(m=b+1\) and \(T=t+1\).  Lean represents the coefficient table by ten
pairwise distinct constructors:
\[
\begin{array}{c|c}
b<t&(m,T)=(\mathrm E,\mathrm E),(\mathrm E,\mathrm O),
       (\mathrm O,\mathrm E),(\mathrm O,\mathrm O),\\
b=t&T=\mathrm E,T=\mathrm O,\\
b>t&(m,T)=(\mathrm E,\mathrm E),(\mathrm E,\mathrm O),
       (\mathrm O,\mathrm E),(\mathrm O,\mathrm O).
\end{array}
\]
The classifier first compares \(b\) and \(t\), then dispatches on the
conductor parities.  It is exhaustive, and injectivity of a numeric code
proves that its constructors are disjoint.  Its character-parity function is
exactly
\[
\begin{array}{c|cccc}
 &e_K&e_\tau&e_\chi&e_{\tau\chi}\\ \hline
b<t&\epsilon(m)&\epsilon(T)&\epsilon(m)&\epsilon(T)\\
b=t&\epsilon(T)&\epsilon(T)&\epsilon(T)&\epsilon(T)\\
b>t&\epsilon(T)&\epsilon(T)&\epsilon(m)&\epsilon(m).
\end{array}
\]

The exact all-even predicate has precisely three constructors,
\[
 \texttt{belowMEvenTEven},\qquad
 \texttt{boundaryEven},\qquad
 \texttt{aboveMEvenTEven}.
\]
These, and only these, have parity vector
\((e_K,e_\tau,e_\chi,e_{\tau\chi})=(0,0,0,0)\).  The complementary seven
rows carry an actual positive-polar source:
\[
\begin{array}{c|c|c}
\text{row}&(e_K,e_\tau,e_\chi,e_{\tau\chi})&
  \text{retained lower source}\ \\ \hline
\texttt{belowMEvenTOdd}&(0,1,0,1)&\tau\\
\texttt{belowMOddTEven}&(1,0,1,0)&\chi_F\\
\texttt{belowMOddTOdd}&(1,1,1,1)&\tau,\chi_F\\
\texttt{boundaryOdd}&(1,1,1,1)&\tau,\chi_F\\
\texttt{aboveMEvenTOdd}&(1,1,0,0)&\tau\\
\texttt{aboveMOddTEven}&(0,0,1,1)&\chi_F\\
\texttt{aboveMOddTOdd}&(1,1,1,1)&\tau,\chi_F.
\end{array}
\]
Here ``source'' means one of the actual quotient-derived lower critical
functions already retained by \texttt{WildQuadraticCoefficientInputs}, not
an arbitrary Hasse function.  Lean proves the all-even/positive-polar
dichotomy, its disjointness, and the equivalence between positive-polar
provenance and failure of the exact all-even predicate.

The genuinely common coefficient data are factored into
\texttt{WildQuadraticCommonCoefficientView}.  This record is independent of
any characteristic-two refinement and retains the actual conductor row,
common correction, selected stationary pairs and their exact ratios, local
character identifications, and the common correction coordinate.  Its
\(z_0\)-unit, \(z_1\)-unit, and residue-field equivalence are the same
objects exposed by the refined view.  The old
\texttt{WildQuadraticCoefficientView} keeps its public signature and
forwards these fields through \texttt{toCommon}.

An exact \texttt{WildQuadraticAllEvenCoefficientView} consists of this
q-free common view together with an \texttt{IsAllEven} witness.  Its three
constructors correspond exactly to the three displayed manuscript rows and
none accepts a refinement.  Its coefficient data are the literal constant
rows, and \texttt{WildQuadraticAllEvenCoefficientTable} exposes the row,
stationary-pair, character, correction, unit, and residue-coordinate
provenance without introducing \(\mathfrak q\).

On the complementary branch,
\texttt{WildQuadraticCoefficientInputs.PositivePolarSource} retains the
actual lower quotient-derived Hasse function from which normalization is
allowed.  A \texttt{WildQuadraticNormalizedRefinementSource} contains a
normalized refinement together with its proof of provenance from that
specific source.  Thus every refined row has a source-tied \(\mathfrak q\),
whereas no refinement or Hasse source can occur in an all-even view.

An even entry has a separate constant constructor whose residual function is
literally \(1\).  An odd entry is
\[
 h_\lambda(z)=\mathfrak q(z)\psi _0(\lambda z),
\]
where \(\mathfrak q\) is the chosen normalized characteristic-two
refinement.  Thus no affine or polar coefficient is assigned to an even
entry.  For an odd entry the intrinsic polar coefficient is the positive
coefficient \(1\).  The intrinsic affine coefficients of
\((\chi_K,\tau,\chi_F,\tau\chi_F)\) in the ten rows are
\[
\begin{array}{c|c}
b<t,\ (\mathrm E,\mathrm E)&(-,-,-,-)\\
b<t,\ (\mathrm E,\mathrm O)&(-,\gamma _0,-,\gamma _0)\\
b<t,\ (\mathrm O,\mathrm E)&(\gamma',-,\gamma,-)\\
b<t,\ (\mathrm O,\mathrm O)&(\gamma',\gamma _0,\gamma,\gamma _0)\\
b=t,\ \mathrm E&(-,-,-,-)\\
b=t,\ \mathrm O&(\gamma',\gamma _0,\gamma,\gamma _1)\\
b>t,\ (\mathrm E,\mathrm E)&(-,-,-,-)\\
b>t,\ (\mathrm E,\mathrm O)&(\gamma',\gamma _0,-,-)\\
b>t,\ (\mathrm O,\mathrm E)&(-,-,\gamma,\gamma)\\
b>t,\ (\mathrm O,\mathrm O)&(\gamma',\gamma _0,\gamma,\gamma).
\end{array}
\]
Here a dash means the literal constant factor.  The exact gamma table is
\[
\begin{array}{c|c}
b<t&(\gamma')^2=\gamma,\qquad\gamma _1=\gamma _0,\\
b>t&(\gamma')^2=1+\gamma _0,\qquad\gamma _1=\gamma,\\
b=t&\displaystyle
 \gamma _1=\frac{\gamma _0+\rho\gamma+r}{1+\rho},\qquad
 (\gamma')^2=\frac{\gamma+\rho\gamma _0+\rho+r}{1+\rho}.
\end{array}
\]
At the odd boundary Lean retains
\(\rho\ne0\), \(r^2=\rho\), and \(1+\rho\ne0\); no denominator is
suppressed.  In the constructed table \(r\) and every \(\gamma'\) are the
unique characteristic-two square roots supplied by
\texttt{QuadraticRefinement}, rather than additional choices.

In the formula-facing constructed table, the lower odd coefficients are not
independently supplied parameters.  A quotient-derived critical witness
retains the exact local characters
\(\chi,\psi\), the equality \(a(\chi)=2d+1\), the hypothesis
\(1<a(\chi)\), an admissible denominator \(\Gamma\), the critical coordinate
\(\delta\) with its exact order, an actual stationary numerator
representative \(\beta\), and the proof that \(\beta\) represents the required
stationary quotient class.  It also retains the exact
\texttt{CriticalPolarCoordinate} equality saying that the polar coefficient
in this coordinate is \(+1\).  Lean then uses uniqueness in
\texttt{QuadraticRefinement} to extract the affine coefficient of this actual
function.  A ten-constructor input type includes precisely the \(\tau\) and
\(\chi_F\) witnesses that occur in each odd row and constructs all other
coefficients from them.  Matching predicates identify their denominators and
numerators with the selected stationary pairs and identify their characters
with the actual local data.  At the odd boundary they additionally retain an
explicit unit lift \(\widetilde\rho\), its residue equality
\(\overline{\widetilde\rho}=\rho\), and the exact relation
\(c_\chi=c_\tau/\widetilde\rho\).  Thus the common coordinate is
choice-relative and no lift is made canonical.

The common residue coordinate uses the scales
\[
 (s_K,s_\tau,s_\chi,s_{\tau\chi})=
 \begin{cases}
 (1,1,1,1),&b\ne t,\\
 (1+r,1,\rho,1+\rho),&b=t\text{ and }T\text{ is odd}.
 \end{cases}
\]
For every odd factor, \texttt{CriticalPolarCoordinate.scale} computes the
transported polar and affine coefficients as
\[
 (A_\theta,B_\theta)=(s_\theta^2,s_\theta\lambda_\theta),
\]
with the positive polar sign.  In particular, at the odd boundary these are
\[
\begin{array}{c|c}
\chi_K&((1+r)^2,(1+r)\gamma'),\\
\tau&(1,\gamma _0),\\
\chi_F&(\rho^2,\rho\gamma),\\
\tau\chi_F&((1+\rho)^2,\gamma _0+\rho\gamma+r).
\end{array}
\]
The raw upper affine coefficient, before Frobenius normalization, is
\(\gamma\) below the break, \(1+\gamma _0\) above the break, and
\[
 L=\frac{\gamma+\rho\gamma _0+\rho+r}{1+\rho}
\]
at the odd boundary.  The raw boundary scale is \(1+\rho\), while the
normalized upper scale is \(1+r\).  Hence the raw common polar and affine
coefficients are \(((1+\rho)^2,
\gamma+\rho\gamma _0+\rho+r)\), and Frobenius transport proves
\[
 h_L((1+\rho)z^2)=h_{\gamma'}((1+r)z).
\]

The finite coefficient data define and prove the exact residue transport
\[
 (p(z),q(z))=
 \begin{cases}
 (z,0),&b<t,\\
 (z,\rho z),&b=t\text{ in the odd row},\\
 (0,z),&b>t,
 \end{cases}
\]
with the boundary-even second coordinate immaterial.  In these coordinates
the raw upper function is
\[
 \Phi_{\chi_K}^{\rm raw}(z)=
 \Phi_{\chi_F}(p(z))\Phi_\tau(q(z))
 \psi _0(q(z))^{e_\tau}.
\]
The boundary lower product and raw upper product are identified with the
normalized refinements from \texttt{QuadraticRefinement}, and every one of the
ten finite constructors satisfies
\(\Phi_{\chi_K}^{\rm raw}(z^2)=\Phi_{\chi_K}(z)\).

The corresponding whole-function identities for the actual quotient-derived
functions are not stored as hypotheses.  Instead Lean retains primitive,
phase-preserving coordinate certificates.  For the lower product these
certificates split exhaustively into the below, boundary, and above ranges.
Off the boundary
they record the dominant quotient witness together with one exact statement
that the complete field-level stationary factor of the stable factor is
\(1\), at the same retained unit.  At the boundary they retain three integral
lifts, their exact reductions to the displayed arguments, and equality of
the resulting local critical units.  Multiplication of the selected
field-level stationary factors then proves the complete lower-product
function.

Residue degree one makes the extension residue map into the canonical
residue-field equivalence used upstairs; Lean proves pointwise that this
equivalence is the extension residue map.  For the upper function the
primitive certificate is supplied with one coordinate point for every
\(z\).  Writing \(w=s_K^{-1}z\), that point retains an integral upper lift
reducing to the canonical image of \(z\), the unit \(1+ux\), its two norm
coordinates, and lower integral lifts reducing exactly to \(p(w^2)\) and
\(q(w^2)\).  In an odd lower row those lifts are tied to the corresponding
quotient witnesses and local units; in an even row the complete field-level
stationary factor at the retained unit is proved to be \(1\).  The exact norm
polynomial identity, the selected stationary ratios, and triviality of the
norm character prove the upper stationary factor as the lower \(\chi_F\)
factor times the inverse \(\tau\) factor.  Positive polar form then proves
\(\Phi_\tau(q)^{-1}=\Phi_\tau(q)\psi_0(q)^{e_\tau}\), yielding the displayed
upper transport.  The coefficient table does not construct the primitive
coordinate points, and neither primitive certificate assumes an equality
between an actual critical function and a table function.

Representative dependence is explicit.  For two actual numerators
\(\beta,\beta'\) proved to represent the same stationary quotient class,
\texttt{CriticalPolarCoordinate} supplies their intrinsic translation
parameter \(a\).  Lean proves both the function identity and, after extracting
the normalized affine coordinates, the exact intrinsic direction
\[
 B'=B+a;
\]
this is \(B'=B+Aa\) with the retained positive normalization \(A=1\).
No canonical stationary representative is introduced.  For a nonzero common
scale \(s\), Lean separately proves
\[
 sB'=sB+s^2(a/s).
\]
Thus \(B_{\rm com}=sB\), \(A_{\rm com}=s^2\), and the common-coordinate
translation parameter is \(a_{\rm com}=a/s\).  Changing to the second
normalized refinement replaces every existing \(\lambda\) by \(\lambda+1\)
without changing its critical function.  Constants remain constants.

The phase quotient preserves the manuscript's numerator and denominator
orientation:
\[
 \frac{H(\gamma')^{e_K}H(\gamma _0)^{e_\tau}}
      {H(\gamma)^{e_\chi}H(\gamma _1)^{e_{\tau\chi}}}.
\]
For the actual phase assembly, the \(\tau\)- and \(\chi_F\)-phases are taken
directly from the retained quotient witnesses and are literally \(1\) when
the corresponding source is absent.  When a derived \(\chi_K\)- or
\(\tau\chi_F\)-entry is odd, its phase has existential quotient-derived
provenance matching the exact selected stationary pair and local characters;
when that entry is even, its value is \(1\).  Lean proves that the product of
the actual \(\chi_K\)- and \(\tau\)-phases divided by the product of the actual
\(\chi_F\)- and \(\tau\chi_F\)-phases is exactly the finite coefficient data's
phase ratio.

The four selected Lamprecht numerator/denominator pairs are also retained
without changing the choices made by the parameter theorems.  In the high
range, with \(c=\alpha\delta_h\), they are
\[
 (\gamma_F,c(n-u)),\quad(\gamma_F/c,1),\quad
 (\gamma_F,cn),\quad(\gamma_F,c(n+1));
\]
the essential quadratic correction unit \(\delta_h\) remains in the
denominator.  In the low range they are
\[
 (\delta/\varepsilon _1,(\alpha/\varepsilon _1)(n-u)),\quad
 (\delta/\alpha,1),\quad
 (\delta/\varepsilon,(\alpha/\varepsilon)n),\quad
 (\delta,\alpha(n+1)).
\]

The low and high packages retain those exact selected representatives, the
actual local data, the quotient-derived odd inputs, the primitive coordinate
certificates just described, and the common-correction record built from the
same data.  These packages take the choice-relative lift, reduction, and
stable-cancellation certificates as explicit proof data.  The transport
theorems on the retained source certificates derive the residual-function
identities.  At the boundary, a linked record keeps the very same
choice-relative lift \(\widetilde\rho\) used by the two lower functions and
proves that its residue, after the canonical extension residue transport,
equals the residue of \texttt{CommonCorrection}'s selected boundary unit.  It
also proves that the residue of the selected \(n\) is \(\rho^2\).  This is an
equality of transported residue classes, not a field-level equality or a
canonical choice of lift.  The formula-facing view assumes \(1<m\), fixes its
row from the actual indices \((m-1,t)\), constructs rather than accepts its
coefficient data, retains the exact three nontrivial stationary-ratio
identities and the opposite-noncancellation proof, and exposes the resulting
units \(z_0,z_1\).
The principal theorem returns this view-level provenance, the primitive
source certificates, the actual phase assembly, the linked boundary
coordinate, and the complete finite coefficient certificate.  Separate low-
and high-package specializations remain indexed by the full original package,
so the selected low record and the full high record containing \(\delta_h\)
are not erased.
No later error, boundary, nonboundary, main, or phase-reduction theorem is
imported or assumed.

\LeanFile{LanglandsFirstMainLemma/Cases/WildQuadratic/CoefficientTable.lean}
\lean{LanglandsFirstMainLemma.WildQuadraticCoefficientRow}
\lean{LanglandsFirstMainLemma.WildQuadraticCoefficientRow.code_injective}
\lean{LanglandsFirstMainLemma.WildQuadraticCoefficientRow.exhaustive}
\lean{LanglandsFirstMainLemma.WildQuadraticCoefficientRow.ofConductors_mParity}
\lean{LanglandsFirstMainLemma.WildQuadraticCoefficientRow.ofConductors_TParity}
\lean{LanglandsFirstMainLemma.WildQuadraticCoefficientRow.IsAllEven}
\lean{LanglandsFirstMainLemma.WildQuadraticCoefficientRow.IsPositivePolar}
\lean{LanglandsFirstMainLemma.WildQuadraticCoefficientRow.isAllEven_or_isPositivePolar}
\lean{LanglandsFirstMainLemma.WildQuadraticCoefficientRow.isAllEven_isPositivePolar_disjoint}
\lean{LanglandsFirstMainLemma.WildQuadraticCoefficientData}
\lean{LanglandsFirstMainLemma.WildQuadraticCoefficientData.coefficient_exists_iff_odd}
\lean{LanglandsFirstMainLemma.WildQuadraticCoefficientData.boundaryOdd_coefficient_table}
\lean{LanglandsFirstMainLemma.WildQuadraticCoefficientData.commonFunction_positivePolar}
\lean{LanglandsFirstMainLemma.WildQuadraticCoefficientData.rawUpper_exact_transport}
\lean{LanglandsFirstMainLemma.WildQuadraticCoefficientData.rawUpper_frobenius}
\lean{LanglandsFirstMainLemma.QuotientDerivedNormalizedCriticalFunction}
\lean{LanglandsFirstMainLemma.QuotientDerivedNormalizedCriticalFunction.function_eq_critical}
\lean{LanglandsFirstMainLemma.QuotientDerivedNormalizedCriticalFunction.toCriticalPolarFunction_inv}
\lean{LanglandsFirstMainLemma.QuotientDerivedNormalizedCriticalFunction.affineCoefficient_replaceRepresentative}
\lean{LanglandsFirstMainLemma.QuotientDerivedNormalizedCriticalFunction.commonAffineCoefficient_replaceRepresentative}
\lean{LanglandsFirstMainLemma.WildQuadraticCoefficientInputs}
\lean{LanglandsFirstMainLemma.WildQuadraticCoefficientInputs.PositivePolarSource}
\lean{LanglandsFirstMainLemma.WildQuadraticCoefficientInputs.PositivePolarSource.toHasse}
\lean{LanglandsFirstMainLemma.WildQuadraticCoefficientInputs.isAllEven_or_positivePolarSource}
\lean{LanglandsFirstMainLemma.WildQuadraticCoefficientInputs.isAllEven_positivePolarSource_disjoint}
\lean{LanglandsFirstMainLemma.WildQuadraticCoefficientInputs.positivePolarSource_iff_not_isAllEven}
\lean{LanglandsFirstMainLemma.WildQuadraticNormalizedRefinementSource}
\lean{LanglandsFirstMainLemma.WildQuadraticNormalizedRefinementSource.ofPositivePolarSource}
\lean{LanglandsFirstMainLemma.WildQuadraticNormalizedRefinementSource.exists_of_positivePolarSource}
\lean{LanglandsFirstMainLemma.WildQuadraticCoefficientInputs.toCoefficientData}
\lean{LanglandsFirstMainLemma.WildQuadraticCoefficientInputs.WildQuadraticBoundaryCommonCoordinate}
\lean{LanglandsFirstMainLemma.WildQuadraticCoefficientInputs.WildQuadraticBoundaryCommonCoordinate.chiF_common_polarCoefficient}
\lean{LanglandsFirstMainLemma.ActualProductCoordinateCertificate.transport}
\lean{LanglandsFirstMainLemma.ActualUpperCoordinateCertificate.transport}
\lean{LanglandsFirstMainLemma.wildQuadraticResidueEquiv_apply}
\lean{LanglandsFirstMainLemma.WildQuadraticBoundaryCorrectionCoordinate}
\lean{LanglandsFirstMainLemma.WildQuadraticCoefficientInputs.MatchesBoundaryCorrectionCoordinate}
\lean{LanglandsFirstMainLemma.WildQuadraticActualPhaseAssembly}
\lean{LanglandsFirstMainLemma.wildQuadraticActualPhaseRatio_exact}
\lean{LanglandsFirstMainLemma.lowWildQuadraticSelectedStationaryPairs_exact}
\lean{LanglandsFirstMainLemma.highWildQuadraticSelectedStationaryPairs_exact}
\lean{LanglandsFirstMainLemma.WildQuadraticCommonCoefficientView}
\lean{LanglandsFirstMainLemma.WildQuadraticCommonCoefficientView.z0Unit}
\lean{LanglandsFirstMainLemma.WildQuadraticCommonCoefficientView.z1Unit}
\lean{LanglandsFirstMainLemma.WildQuadraticCommonCoefficientView.residueEquiv}
\lean{LanglandsFirstMainLemma.WildQuadraticAllEvenCoefficientView}
\lean{LanglandsFirstMainLemma.WildQuadraticAllEvenCoefficientView.ofBelowMEvenTEven}
\lean{LanglandsFirstMainLemma.WildQuadraticAllEvenCoefficientView.ofBoundaryEven}
\lean{LanglandsFirstMainLemma.WildQuadraticAllEvenCoefficientView.ofAboveMEvenTEven}
\lean{LanglandsFirstMainLemma.WildQuadraticCoefficientView}
\lean{LanglandsFirstMainLemma.WildQuadraticCoefficientView.toCommon}
\lean{LanglandsFirstMainLemma.WildQuadraticCoefficientView.toAllEven}
\lean{LanglandsFirstMainLemma.WildQuadraticLowCoefficientPackage}
\lean{LanglandsFirstMainLemma.WildQuadraticHighCoefficientPackage}
\lean{LanglandsFirstMainLemma.WildQuadraticAllEvenCoefficientCertificate}
\lean{LanglandsFirstMainLemma.wildQuadraticAllEvenCoefficientCertificate}
\lean{LanglandsFirstMainLemma.WildQuadraticAllEvenCoefficientTable}
\lean{LanglandsFirstMainLemma.WildQuadraticCoefficientTable}
\lean{LanglandsFirstMainLemma.WildQuadraticLowCoefficientTable}
\lean{LanglandsFirstMainLemma.WildQuadraticHighCoefficientTable}
\leanok
\SourceText{Lemma lem:critical-polar-coordinate; Propositions prop:quadratic-critical-functions and prop:quadratic-gamma-table; the surrounding wild-quadratic coefficient and parity tables.}
\uses{mod:cases-wildquadratic-commoncorrection,mod:cases-wildquadratic-quadraticrefinement,mod:lamprecht-criticalpolarcoordinate,mod:parameters-minimalorbitstationary}
\FormalizationContract
\end{definition}
